\chapter{기호}
\label{sec:notation}

뉴런의 유형은 주 신경전달물질 영어 이름의 첫 네 글자로 표기한다. \nt{GLUT}는 글루탐산, \nt{GABA}는 감마-아미노뷰티르산, \nt{GLYC}는 글리신, \nt{ACET}는 아세틸콜린, \nt{DOPA}는 도파민, \nt{SERO}는 세로토닌, \nt{HIST}는 히스타민, \nt{NORA}는 노르아드레날린, \nt{ADRE}는 아드레날린, \nt{ASPA}는 아스파르트산이다(표~\ref{tab:types}).

문자는 양보다 적어서, 장에 따라 같은 문자가 다른 것을 뜻하기도 한다. 아래 표에서는 뜻마다 행을 따로 두었다. 오른쪽 열은 그 기호를 도입한 식이다.

\begin{center}
\footnotesize
\setlength{\tabcolsep}{4pt}
\renewcommand{\arraystretch}{1.12}
\begin{longtable}{>{\raggedright}p{2.0cm}>{\raggedright}p{9.6cm}>{\raggedright\arraybackslash}p{2.4cm}}
\toprule
기호 & 뜻 & 도입한 곳 \\
\midrule
\endfirsthead
\toprule
기호 & 뜻 & 도입한 곳 \\
\midrule
\endhead
\bottomrule
\endfoot
$a$ & 선형 전달 특성의 기울기, $f(u)=au$ & \eqref{eq:feedback} \\
$a$ & \nt{ACET} 수준: $0$은 읽기, $1$은 쓰기 & \eqref{eq:acetnet} \\
$a_i$, $a$ & 집단의 피로: 리듬 발생기에서; 서파 수면의 시소에서 & \eqref{eq:halfcentre}, \eqref{eq:updown} \\
$a_S$, $a_P$ & 가속과 제동에 대한 심장 박동의 민감도 & \eqref{eq:gasbrake} \\
$A(r)$, $A_0$ & 빈도 $r$일 때 스파이크에 대한 시냅스 응답; 쉬고 난 시냅스의 응답 & \eqref{eq:depression} \\
$A_1$ & 전달물질 한 꾸러미에 대한 수신자의 응답 & \eqref{eq:quantal} \\
$A$ & 수상돌기를 포함한 뉴런의 막 면적 & \eqref{eq:dendriteload} \\
$A_f$, $A_s$ & 다음 움직임까지 보존되는 보정의 비율, 빠른 과정과 느린 과정 & \eqref{eq:tworate} \\
$b_i$ & 페이스메이커 자체의 유입 & \eqref{eq:seroneuron} \\
$b$ & 리듬 발생기나 수면 시소에서 활동이 집단을 얼마나 빨리 지치게 하는가 & \eqref{eq:halfcentre}, \eqref{eq:updown} \\
$b$ & 디지털화 표본 하나당 비트 수 & \eqref{eq:probedata} \\
$B_f$, $B_s$ & 보정이 오차로부터 배우는 세기, 빠른 과정과 느린 과정 & \eqref{eq:tworate} \\
$c$ & 부피 속 전달물질의 농도 & \eqref{eq:seroneuron} \\
$c$ & 이웃 뉴런 잡음의 상관계수 & \eqref{eq:correlated} \\
$c$ & 스파이크 수 차이가 라켓을 움직이는 계수 & \eqref{eq:dishreadout} \\
$c_f$ & $C$ 뉴런의 응답: 모든 위치에 걸친 $s_{f,p}$의 최댓값 & \eqref{eq:scstage} \\
$c_m$ & 막의 비정전용량 & \eqref{eq:dendriteload} \\
$C$ & 노력의 대가: 시간 $\tau_a$ 동안의 행동은 $C/\tau_a$가 든다 & \eqref{eq:vigor} \\
$d'$ & 두 입력의 변별도: 차이를 잡음으로 나눈 것 & \eqref{eq:gainsnr} \\
$d$, $d_j$ & 신호 경로의 지연 & \eqref{eq:glutneuron}, \eqref{eq:vstar} \\
$d$ & 근육과의 되먹임 루프의 지연 & \eqref{eq:delaystable} \\
$D$ & 확산 계수 & \eqref{eq:diffusionlength} \\
$D$ & \nt{GLUT}--\nt{GABA} 네트워크 발화율 식의 분모 & \eqref{eq:isn} \\
$D$ & 지연된 보상을 기다리는 시간 & \eqref{eq:patience} \\
$e$, $e_{ij}$ & 시냅스의 흔적 & \eqref{eq:threefactor} \\
$e$, $e_i$ & 오차 전선으로 오는 출력 또는 움직임의 오차 & \eqref{eq:lms}, \eqref{eq:adaptivefilter}, \eqref{eq:tworate} \\
$E_{\text{spk}}$ & 시냅스의 일을 포함한 스파이크 하나의 에너지 & \eqref{eq:energy} \\
$E$ & \nt{GLUT} 집단의 발화율 & \eqref{eq:ei} \\
$E_E$ & 흥분성 시냅스의 평형 전위 & \eqref{eq:shunt} \\
$f$, $F$ & 전달 특성: 입력의 함수로서의 발화율 & \eqref{eq:ratenet}, \eqref{eq:gain} \\
$f$, $f_0$ & 심박수; 페이스메이커 고유의 리듬 & \eqref{eq:gasbrake} \\
$f_s$ & 전극 하나의 디지털화 빈도 & \eqref{eq:probedata} \\
$F(t)$, $F_1$ & 근육의 힘; 연축 한 번의 힘 & \eqref{eq:forcesum} \\
$F_1$, $F_2$ & 관절을 반대 방향으로 당기는 두 근육의 힘 & \eqref{eq:antagonists} \\
$g$ & 농도에 대한 수용체의 민감도 & \eqref{eq:seroneuron} \\
$g_L$, $g_E$, $g_I$ & 누설, 흥분, 억제의 전도도 & \eqref{eq:shunt} \\
$g$ & \nt{NORA}가 정하는 이득 & \eqref{eq:gain}, \eqref{eq:machine} \\
$g$ & 기억 네트워크에서 \nt{GABA}에 의한 전체 억제 & \eqref{eq:acetnet} \\
$g$ & 위에서 정하는 통증 관문의 이득 & \eqref{eq:gate} \\
$g$ & 호르몬 연쇄 한 단계의 이득 & \eqref{eq:cascade} \\
$g_j$ & 적응 필터 입력에 있는, 고유 시간 상수를 가진 필터 $j$ & \eqref{eq:adaptivefilter} \\
$G$ & 되먹임 루프의 이득 & \eqref{eq:feedback} \\
$G$ & 지연 있는 루프의 보정 속도 & \eqref{eq:delaystable} \\
$h$, $h_I$ & \nt{GLUT} 집단과 \nt{GABA} 집단의 외부 입력 & \eqref{eq:ei}, \eqref{eq:isn} \\
$h(t)$ & 근육의 단일 연축 & \eqref{eq:twitch} \\
$H(t)$ & 수면 압력의 위쪽 문턱값: 여기에 이르면 기계가 잠든다 & \eqref{eq:sleeppressure} \\
$I$, $I_i$ & 뉴런이나 집단으로의 외부 유입 & \eqref{eq:ratenet}, \eqref{eq:wta}, \eqref{eq:updown} \\
$I$ & \nt{GABA} 집단의 발화율 & \eqref{eq:ei} \\
$I$ & 자극의 밝기 또는 세기 & \eqref{eq:photonnoise}, \eqref{eq:nakarushton} \\
$I$ & 막을 충전하는 입력 전류 & \eqref{eq:dendriteload} \\
$k$ & 버스트 안의 스파이크 수 & \eqref{eq:burst} \\
$k$ & 증가분이 잡음의 몇 배를 넘어야 하는가 & \eqref{eq:photonnoise} \\
$k$ & 호르몬 연쇄의 되먹임 계수 & \eqref{eq:cascade} \\
$k_s$ & 통증이 민감화를 끌어올리는 세기 & \eqref{eq:sensitization} \\
$K$ & 결과의 특징 수 & \eqref{eq:vectortd} \\
$K$ & 관절의 강성 & \eqref{eq:antagonists} \\
$K$ & 호르몬 연쇄의 루프 이득, $K=g^3k$ & \eqref{eq:cascade}, \eqref{eq:cascadestable} \\
$K_\varphi$ & 외부 신호에 대한 일주기 발생기의 동조 세기 & \eqref{eq:pll} \\
$L$ & 지연선의 길이 & \eqref{eq:itd} \\
$L$ & 통증 센서에서 오는 전선의 길이 & \eqref{eq:paintime} \\
$L(t)$ & 수면 압력의 아래쪽 문턱값: 여기에 이르면 기계가 깬다 & \eqref{eq:sleeppressure} \\
$M$ & 예고 신호와 보상 사이 인과 관계의 척도 & \eqref{eq:retro} \\
$M_k$ & 특징 $k$의 기댓값 & \eqref{eq:vectortd} \\
$M_{ik}$ & \nt{GABA} 뉴런 $k$로부터의 억제 세기 & \eqref{eq:machine} \\
$M$ & 관절의 회전력 & \eqref{eq:antagonists} \\
$M$ & 혼합기의 입력선 수 & \eqref{eq:cover} \\
$n$, $\bar n$ & 창 안의 스파이크 수; 그 평균 & \eqref{eq:rateerror} \\
$n$ & 문턱 소자의 입력 수 & \eqref{eq:cover} \\
$n_\uparrow$, $n_\downarrow$ & 창 안에서 `위' 영역과 `아래' 영역의 스파이크 수 & \eqref{eq:dishreadout} \\
$N$ & 뉴런 수 & \eqref{eq:correlated}, \eqref{eq:energy} \\
$N$ & 프로브의 전극 수 & \eqref{eq:probedata} \\
$N_s$ & 두 뉴런 사이의 방출 지점 수 & \eqref{eq:quantal} \\
$p$ & 스파이크당 전달물질 방출 확률 & \eqref{eq:burst} \\
$p$ & 기억 네트워크에서 활성 뉴런의 비율 & \eqref{eq:acetnet} \\
$p$ & 보상하도록 학습하는 섭동 & \eqref{eq:tworate} \\
$P$ & 신호에 쓰는 전력 & \eqref{eq:energy}, \eqref{eq:dishpower} \\
$P$ & 저장된 추정값의 불확실성 & \eqref{eq:kalman} \\
$P$ & `제동'의 활동: 심장으로 가는 \nt{ACET} 전선 & \eqref{eq:gasbrake} \\
$P_{\max}$ & 문턱 소자가 두 클래스로 나눌 수 있는 무작위 패턴의 수 & \eqref{eq:cover} \\
$P_{\text{miss}}$ & 미스 확률 & \eqref{eq:dishdensity} \\
$q$ & 세상이 변하는 속도 & \eqref{eq:kalman} \\
$q_k$ & \nt{GABA} 뉴런 $k$의 발화율 & \eqref{eq:machine} \\
$q_0$ & 억제성 중개자의 배경 활동 & \eqref{eq:gate} \\
$q$ & 통증 관문에서 억제성 중개자의 발화율 & \eqref{eq:gate} \\
$q$ & 다음 운동 단위가 앞의 것보다 몇 배 센가 & \eqref{eq:sizeprinciple} \\
$r$, $r_i$ & 뉴런의 스파이크 발화율 & \eqref{eq:ratenet} \\
$r$, $r_t$ & 보상(보상의 흐름) & \eqref{eq:rpe}, \eqref{eq:ctd} \\
$\mathbf r(t)$ & 저수지의 응답 벡터 & \eqref{eq:reservoir} \\
$R$, $\bar R$ & 받은 보상; 그 기댓값 또는 단위 시간당 평균 & \eqref{eq:perturbation}, \eqref{eq:vigor} \\
$R$ & 필터 입력 잡음의 분산 & \eqref{eq:kalman} \\
$R$, $r_0$ & 지연된 보상과 즉각적인 보상 & \eqref{eq:patience} \\
$R_{\max}$ & 최대 전송 속도, bit/s & \eqref{eq:spikeentropy} \\
$r_{\max}$ & 최대 스파이크 발화율 & \eqref{eq:gain}, \eqref{eq:nakarushton} \\
$R_{\text{raw}}$, $R_{\text{spk}}$ & 프로브의 데이터 흐름: 원시 데이터와 스파이크만, bit/s & \eqref{eq:probedata} \\
$s_j$ & 뉴런 $j$ 출력의 부호 & \eqref{eq:dalesign} \\
$s_i$ & 수신자 수용체의 부호 & \eqref{eq:seroneuron} \\
$s$, $\hat s$ & 세상의 상태; 그 추정값 & \eqref{eq:rpe}, \eqref{eq:kalman} \\
$s_{f,p}$ & 위치 $p$에서 특징 $f$에 대한 $S$ 뉴런의 응답 & \eqref{eq:scstage} \\
$S$ & `가속'의 활동: 심장으로 가는 \nt{NORA} 전선 & \eqref{eq:gasbrake} \\
$S$ & 수면 압력 & \eqref{eq:sleeppressure} \\
$t_1$ & 첫 스파이크의 시각 & \eqref{eq:latency} \\
$t_k$ & $k$번째 스파이크의 시각 & \eqref{eq:forcesum} \\
$t_{f,i}$ & 특징 $f$의 견본 & \eqref{eq:scstage} \\
$T$ & 스파이크를 세는 창; 광자를 모으는 시간 & \eqref{eq:rateerror}, \eqref{eq:photonnoise} \\
$T_c$ & 근육 연축의 상승 시간 & \eqref{eq:twitch} \\
$T_{\text{burst}}$ & 배양의 일제 발화 간격 & \eqref{eq:burstinterval} \\
$u$ & 뉴런의 총입력 & \eqref{eq:gain} \\
$u$, $u(t)$ & 뇌의 명령: 적응 필터의 입력에서; 호르몬 연쇄로 & \eqref{eq:adaptivefilter}, \eqref{eq:cascade} \\
$u(t)$ & 저수지의 입력 & \eqref{eq:reservoir} \\
$U$ & 일정한 입력의 수준 & \eqref{eq:latency} \\
$U$ & 스파이크마다 방출되는 전달물질 저장량의 비율 & \eqref{eq:depression} \\
$v$ & 전선 속 신호의 속도; 점이 움직이는 속도 & \eqref{eq:itd}, \eqref{eq:vstar} \\
$V$ & 막 전압 & \eqref{eq:glutneuron}, \eqref{eq:shunt} \\
$V$ & 미래 보상의 기댓값 & \eqref{eq:rpe}, \eqref{eq:ctd} \\
$w$, $w_{ij}$, $W_{ij}$ & 연결 가중치 & \eqref{eq:glutneuron}, \eqref{eq:ratenet}, \eqref{eq:acetnet}, \eqref{eq:lms}, \eqref{eq:adaptivefilter} \\
$\mathbf w_k$ & $C$ 뉴런 $k$의 입력 가중치 & \eqref{eq:traceinvariance} \\
$w_{q\mu}$, $w_{q\nu}$, $w_{z\mu}$ & 통증 관문의 연결 가중치 & \eqref{eq:gate} \\
$W_{\text{out}}$ & 저수지 선형 판독의 가중치 & \eqref{eq:reservoir} \\
$x$, $y$ & 논리 입력과 출력 & \eqref{eq:dualrail}, \eqref{eq:veto} \\
$x$, $y$ & 송신자와 수신자의 활동 & \eqref{eq:hebb} \\
$x$ & 지연선 위 검출기의 위치 & \eqref{eq:itd} \\
$x_i$ & 감각 기관에서 오는 입력 & \eqref{eq:acetnet} \\
$x_p$ & 위치 $p$의 입력 & \eqref{eq:scstage} \\
$\mathbf x$ & $S$ 뉴런의 출력, $C$ 뉴런의 입력 & \eqref{eq:traceinvariance} \\
$x_j$ & 적응 필터의 입력 $j$: 필터 $g_j$를 거친 명령 & \eqref{eq:lms}, \eqref{eq:adaptivefilter} \\
$x_f$, $x_s$ & 빠른 과정과 느린 과정의 보정 & \eqref{eq:tworate} \\
$x_1$, $x_2$, $x_3$ & 호르몬 연쇄 세 단계의 수준; $x_3$은 최종 호르몬 & \eqref{eq:cascade} \\
$x^*$ & 새 사태가 시작되는 전달물질 저장량의 회복 비율 & \eqref{eq:burstinterval} \\
$y$ & 필터의 잡음 섞인 입력 & \eqref{eq:kalman} \\
$y$, $\hat y$ & 센서 신호; 적응 필터의 예측 & \eqref{eq:adaptivefilter} \\
$y_k$, $\bar y_k$ & $C$ 뉴런 $k$의 활동; 그 흔적 & \eqref{eq:traceinvariance} \\
$y(t)$ & 저수지 판독의 출력 & \eqref{eq:reservoir} \\
$z$ & 통증 관문 출력의 발화율 & \eqref{eq:gate} \\
\midrule
$\alpha_i^\pm$ & 양의 오차와 음의 오차의 가중치 & \eqref{eq:asymmetric} \\
$\beta$ & 상호 억제의 세기 & \eqref{eq:wta} \\
$\beta$ & 통증 관문 출력에 대한 억제의 세기 & \eqref{eq:gate} \\
$\gamma$ & 할인 계수 & \eqref{eq:rpe} \\
$\delta(\cdot)$ & 델타 함수: 순간적인 도약 & \eqref{eq:glutneuron} \\
$\delta$, $\delta_t$ & 보상 예측 오차 & \eqref{eq:rpe}, \eqref{eq:ctd} \\
$\delta t$ & 소리가 두 귀에 도착하는 시간 차 & \eqref{eq:itd} \\
$\Delta$, $\Delta_{\max}$ & 스파이크 간격; 동시성 창 & \eqref{eq:window} \\
$\Delta t$ & 수신자가 시간을 구별하는 정밀도 & \eqref{eq:spikeentropy} \\
$\Delta F$ & 다음 운동 단위가 더하는 힘 & \eqref{eq:sizeprinciple} \\
$\Delta I$ & 구별할 수 있는 밝기 증가분 & \eqref{eq:photonnoise} \\
$\Delta p$ & 창 하나 동안의 라켓 이동 & \eqref{eq:dishreadout} \\
$\Delta u$ & 두 집단의 입력 차이 & \eqref{eq:gainsnr} \\
$\Delta V$ & 막 전압을 얼마나 올려야 하는가 & \eqref{eq:dendriteload} \\
$\Delta x$ & 이웃 센서 사이의 거리 & \eqref{eq:vstar} \\
$\varepsilon$ & 입력의 상대적 차이 & \eqref{eq:wtagain} \\
$\eta$ & 학습률 & \eqref{eq:plasticity}, \eqref{eq:lms}, \eqref{eq:traceinvariance} \\
$\eta$ & 증폭기 뒤 네트워크의 잡음 & \eqref{eq:softmax} \\
$\mu$ & 통증 관문으로 가는 촉각 입력 & \eqref{eq:gate} \\
$\nu$ & 통증 입력 & \eqref{eq:gate} \\
$\theta$ & 발화 문턱값 & \eqref{eq:window}, \eqref{eq:scstage} \\
$\kappa$ & 스파이크 하나당 전달물질의 양 & \eqref{eq:seroneuron} \\
$\kappa_i$ & 양의 오차와 음의 오차 가중치의 비 & \eqref{eq:expectile} \\
$\kappa_m$ & 근육의 힘과 강성 사이의 관계 & \eqref{eq:antagonists} \\
$\ell$ & 전달물질이 퍼지는 거리 & \eqref{eq:diffusionlength} \\
$\ell$ & 근육이 관절을 당기는 팔의 길이 & \eqref{eq:antagonists} \\
$\xi$ & 뉴런 출력의 무작위 흔들림 & \eqref{eq:perturbation} \\
$\rho(\boldsymbol\psi)$ & 설정 $\boldsymbol\psi$에 머문 시간의 비율 & \eqref{eq:dishdensity} \\
$\sigma$ & 산포, 잡음의 크기 & \eqref{eq:rateerror}, \eqref{eq:softmax} \\
$\sigma$ & 응답이 최댓값의 절반이 되는 밝기 & \eqref{eq:nakarushton} \\
$\sigma_{\text{pre}}$, $\sigma_{\text{post}}$ & 증폭기 앞과 뒤의 잡음 & \eqref{eq:gainsnr} \\
$\tau$ & 막의 시간 상수 & \eqref{eq:glutneuron} \\
$\tau$ & 호르몬 연쇄 한 단계의 시간 상수 & \eqref{eq:cascade} \\
$\tau_{\text{eff}}$ & 스스로를 흥분시키는 집단의 시간 상수 & \eqref{eq:perfectint} \\
$\tau_c$ & 부피 속 전달물질의 수명 & \eqref{eq:seroneuron} \\
$\tau_E$, $\tau_I$ & \nt{GLUT} 집단과 \nt{GABA} 집단의 시간 상수 & \eqref{eq:ei} \\
$\tau_e$ & 흔적의 수명 & \eqref{eq:threefactor} \\
$\tau_{\text{tr}}$ & $C$ 뉴런 활동 흔적의 수명 & \eqref{eq:traceinvariance} \\
$\tau_s$ & 손상 뒤 민감도가 돌아오는 시간 & \eqref{eq:sensitization} \\
$\tau_r$, $\tau_d$ & 깨어 있을 때 수면 압력이 쌓이는 시간과 잠잘 때 풀리는 시간 & \eqref{eq:sleeppressure} \\
$\tau_{\text{rec}}$ & 전달물질 저장량의 회복 시간 & \eqref{eq:depression} \\
$\tau_\gamma$ & 계획의 시간 지평 & \eqref{eq:ctd} \\
$\tau_v$ & 기댓값이 갱신되는 속도 & \eqref{eq:ramp} \\
$\tau_a$ & 리듬 발생기나 수면 시소의 피로 시간 & \eqref{eq:halfcentre}, \eqref{eq:updown} \\
$\tau_a^*$ & 행동을 수행하는 최적 시간 & \eqref{eq:vigor} \\
$\Phi$ & 학습 규칙의 우변 & \eqref{eq:plasticity} \\
$\Phi$ & 광자 흐름 & \eqref{eq:photonnoise} \\
$\varphi_k$ & 받은 결과의 특징 $k$ & \eqref{eq:vectortd} \\
$\varphi$ & 일주기 발생기와 외부 신호의 위상 차 & \eqref{eq:pll} \\
$\boldsymbol\psi$ & DishBrain 모델에서 네트워크의 설정 & \eqref{eq:dishdensity} \\
$\omega$, $\omega_0$ & 외부 신호(빛)의 주파수; 일주기 발생기의 고유 주파수 & \eqref{eq:pll} \\
\end{longtable}
\end{center}
